\chapter{Conclusion and Outlook}
\label{cha:conclusion_outlook}

Two-dimensional light-element xenes have emerged as a rapidly developing class of quantum materials.
Their low atomic mass, strong covalent bonding, and structural flexibility give them access to electronic, magnetic, superconducting, topological, and optical phenomena.
Graphene provides the archetypal example, but its zero band gap limits several electronic and optoelectronic applications.
This limitation motivates the exploration of alternative elemental and compound xenes, especially boron-based and beryllium-based systems, where the atomic arrangement can generate a wider range of electronic behaviour.

This thesis investigates three groups of light-element two-dimensional systems using first-principles calculations based on density functional theory.
The calculations determine atomic structures, electronic band structures, densities of states, dynamical stability, dielectric response, and derived optical properties.
For selected systems, the analysis also examines superconductivity, magnetism, and topology.
Across these studies, heterostructure formation, hydrogen functionalization, and reduced dimensionality provide practical routes to connect atomic structure with emergent quantum properties.

Chapter~\ref{cha:project_bilayers} investigates graphene-based van der Waals heterostructures formed from graphene, borophene, and two-dimensional boron carbides, namely \ce{BC3} and \ce{B4C3}.
This study compares the isolated monolayers with the corresponding bilayer heterostructures and determines their electronic and optical properties.
The calculations use the PBE exchange-correlation functional with van der Waals corrections for structural optimization, and the HSE06 hybrid functional for more accurate electronic structures and optical properties.
The optical analysis starts from the frequency-dependent dielectric function \(\varepsilon(\omega)\), and then derives the absorption coefficient, refractive index, extinction coefficient, reflectivity, and energy-loss spectrum.

The isolated monolayers display distinct electronic characters.
Graphene retains its semimetallic band structure, whereas borophene is metallic.
In contrast, monolayer \ce{BC3} and monolayer \ce{B4C3} are indirect-gap semiconductors, with HSE06 band gaps of 1.822 eV and 2.381 eV, respectively.
These contrasting electronic properties provide a basis for heterostructure design.
In the bilayer systems, Graphene-Borophene remains metallic, whereas Graphene-\ce{BC3} and Graphene-\ce{B4C3} are semimetallic.
The graphene-derived linear bands remain near the \(\mathrm{K}\) point, with only small band-gap openings of 20 meV to 50 meV.
Charge-density-difference calculations further show that stacking geometry controls the interfacial charge redistribution between graphene and the boron-containing layer.

The calculated optical properties demonstrate that all three heterostructures absorb in the visible region.
The dielectric response also reveals collective electronic excitations, including \(\pi\)-type and \(\pi\)-\(\sigma\)-type plasmon modes.
Graphene-Borophene gives the strongest absorption and supports low-energy plasmonic features.
Graphene-\ce{B4C3} displays non-negligible off-diagonal components in the dielectric tensor \(\varepsilon_{\alpha\beta}(\omega)\), consistent with anisotropic and possibly chiral optical behaviour.
The Schottky barrier analysis further shows that graphene-boron-carbide heterostructures can provide additional control over charge transport.
We therefore arrive at the result that weak van der Waals coupling can tune electronic and optical properties while preserving key features of the constituent monolayers.

Chapter~\ref{cha:project_boron} extends the discussion from engineered heterostructures to intrinsic quantum behaviour in penta-bipyramid boron phases.
This chapter studies bulk \ce{o-B14}, monolayer \ce{o-B14}, bilayer \ce{o-B14}, and the hydrogen-terminated bilayer.
The edge-sharing pentagonal bipyramids in \ce{o-B14} generate an unusual bonding network, and this network gives boron access to several quantum phases within closely related structures.
The calculations determine the electronic, optical, magnetic, and superconducting properties of these systems.

The two-dimensional \ce{o-B14} structures display a wide range of electronic behaviour.
The bilayer is semiconducting, with an HSE06 band gap of 0.67 eV.
The monolayer shows magnetism, whereas the hydrogen-terminated bilayer becomes superconducting with a critical temperature \(T_\mathrm{c}\) of 24.3 K.
This result demonstrates that hydrogen functionalization can reorganize the electronic structure and activate phonon-mediated superconductivity in a two-dimensional boron system.
The optical calculations also reveal visible-region absorption in selected directions and strong anisotropy in the dielectric response.
For the metallic systems, the real part of the dielectric function changes sign in different directions, producing directional plasmonic modes.
We therefore arrive at the result that the electron-deficient bonding of boron can support semiconductivity, magnetism, superconductivity, and anisotropic plasmonic response within one structural family.

Chapter~\ref{cha:project_beryllene} investigates pristine and hydrogen-functionalized two-dimensional beryllene phases.
Beryllium represents an extreme and relatively unexplored limit among elemental two-dimensional materials.
It is the lightest alkaline-earth metal, and bulk beryllium combines low density, high elastic modulus, high thermal conductivity, and high transparency to X-rays and neutrons.
The two-dimensional form therefore offers a natural platform for examining the effect of reduced dimensionality on the structural, electronic, optical, and topological properties of beryllium.

This chapter considers pristine \(\alpha\)- and \(\beta\)-beryllene, a cubic trilayer beryllene phase, and their hydrogen-functionalized counterparts.
The stability analysis confirms the stability of the previously proposed \(\alpha\)- and \(\beta\)-beryllene phases and identifies a dynamically stable cubic trilayer beryllene phase.
The electronic structures depend strongly on bonding geometry and hydrogen functionalization.
Hydrogen changes the band dispersion, redistributes the electronic states, and modifies the dielectric response.
In particular, hydrogen functionalization drives \(\alpha\)-beryllene into a wide-gap semiconducting state, whereas the hydrogen-functionalized \(\beta\)-based structures retain metallic behaviour.
The optical spectra show strong anisotropy, visible-to-ultraviolet absorption, modified plasmonic behaviour, and direction-dependent dielectric screening.
These results demonstrate that beryllene phases provide an additional light-element platform for tunable optical response.

The topological analysis gives another important result.
Using parity analysis of spin-orbit-coupled electronic states, the calculations identify non-trivial \(\mathbb{Z}_2\) topology in \(\alpha\)-beryllene and cubic trilayer beryllene, whereas \(\beta\)-beryllene remains topologically trivial.
This result shows that topological electronic structure can emerge even in a light-element beryllium system.
Chapter~\ref{cha:project_beryllene} therefore connects structural stability, hydrogen functionalization, optical response, and topology in two-dimensional beryllium phases.

Together, these three studies demonstrate that light-element xenes provide a flexible platform for controlling quantum properties through composition, dimensionality, bonding topology, stacking, and surface functionalization.
The thesis starts from graphene-based heterostructures, proceeds to penta-bipyramid boron phases, and then reaches beryllene as a lighter and less explored elemental limit.
Across these systems, first-principles calculations identify visible optical absorption, anisotropic dielectric response, plasmonic excitations, magnetism, superconductivity, and non-trivial topology.
We therefore arrive at a general picture in which simple light elements can produce complex quantum behaviour when their atomic structures are organized in low-dimensional forms.

Future work can first test these theoretical predictions through experimental synthesis and characterization.
For the boron and beryllium systems, structural verification, vibrational spectroscopy, transport measurements, optical spectroscopy, and electron energy-loss spectroscopy would provide direct tests of the predicted stability, electronic structure, and plasmonic response.
For the superconducting hydrogen-terminated \ce{o-B14} bilayer, measurements of electron transport and magnetic response would clarify the superconducting transition.
For the topological beryllene phases, spectroscopic probes of edge or surface states would help establish the predicted non-trivial topology.

Further theoretical work can also extend the present calculations.
Many-body electronic corrections, excitonic effects, and quantum transport simulations would refine the electronic and optical predictions beyond the independent-particle dielectric response used here.
Strain engineering, chemical doping, hydrogen coverage, substrate effects, moiré superlattices, and external electric fields offer additional routes to tune the band structure, optical response, plasmonic modes, superconductivity, and topology.
These directions connect the material-level predictions in this thesis with possible device-level functions.

The combination of low atomic mass, strong bonding, and tunable low-dimensional electronic structure makes light-element xenes promising candidates for ultrathin superconductors, spintronic materials, plasmonic platforms, and optoelectronic devices.
The results of this thesis therefore provide both specific predictions for boron and beryllium allotropes and a broader route for designing quantum materials from light elements.
